\section{Related Work}

Crack analysis is overwhelmingly detection and
segmentation~\cite{hasan2025review,perez2024assessment}; its datasets annotate
cracks within single images and the large consolidations
merging~\cite{kulkarni2022crackseg9k,benz2024omnicrack30k} inherit that scope,
and multi-view pipelines~\cite{mallios2023multiview} aggregate within one
inspection without instance correspondence. Structural-health monitoring does
compare cracks across time with machinery close to
ours~\cite{parente2022monitoring,harb2025multitemporal,sun2025perspective,nec2019guide},
but identity is always \emph{given by construction} (fixed camera, marked
location, operator at a viewpoint): the subject is one known crack, the question
width or extent, and no dataset supports method comparison on the association
decision itself.

Object re-identification~\cite{qian2024reidtrends} with open-set
identification~\cite{liao2014openset} is the formal setting we adopt; keypoint
matching~\cite{lowe2004sift,rublee2011orb}, learned correspondence
(SuperGlue~\cite{sarlin2020superglue}, LoFTR~\cite{sun2021loftr}), and region
embeddings~\cite{dosovitskiy2021vit,touvron2021deit,radford2021clip,zhou2019osnet,oquab2024dinov2}
apply directly to crack crops and our benchmark measures all of them.
Vehicle-damage assessment already re-identifies damage with OSNet in curated
settings~\cite{maiano2023antifraud}; our results suggest the crop scope is
itself the limiting choice. The move our baseline makes --- align two captures,
then reason about what is present in one and absent in the other --- is change
detection~\cite{sakurada2015changedetection,alcantarilla2018street}; planar
registration is classical~\cite{hartley2004mvg,fischler1987ransac,barath2019magsac}
with photometric alignment (ECC~\cite{evangelidis2008ecc}) as the fallback. Our
contribution is not a new estimator but the observation that the alignment
target and the matching target should be \emph{disjoint}: the homography is
driven by the wall precisely because the crack is what we seek to identify.